%%%%%%%%%%%%%%%%%%%%%%%%%%%%%%%%%%%%%%%%%%%%%%%%%%%%%%%%%%%%
% 14.1 MATHEMATICAL SYMBOLS AND NOTATION
% The Composition of Advanced Mathematics
%%%%%%%%%%%%%%%%%%%%%%%%%%%%%%%%%%%%%%%%%%%%%%%%%%%%%%%%%%%%

\section{Mathematical Symbols and Notation}
\label{sec:mathematical-symbols-notation}

\prerequisite{
Familiarity with the notation introduced throughout the textbook.
}

\learningobjectives{
By the end of this reference section, the reader should be able to:
\begin{itemize}
    \item identify common arithmetic symbols;
    \item identify algebraic symbols;
    \item interpret equality and inequality notation;
    \item identify set-theoretic notation;
    \item identify logical notation;
    \item interpret function notation;
    \item recognize interval notation;
    \item identify number-set symbols;
    \item interpret summation and product notation;
    \item understand factorial and combinatorial notation;
    \item identify calculus notation;
    \item recognize partial derivatives and gradients;
    \item interpret divergence, curl, and Laplacians;
    \item identify vector and matrix notation;
    \item recognize determinants, traces, ranks, and transposes;
    \item interpret eigenvalue notation;
    \item recognize probability notation;
    \item interpret expectation, variance, covariance, and correlation;
    \item recognize statistical notation;
    \item interpret asymptotic notation;
    \item recognize complex-number notation;
    \item identify geometry and trigonometry symbols;
    \item recognize topology notation;
    \item identify abstract algebra notation;
    \item recognize number-theory notation;
    \item identify optimization notation;
    \item interpret differential-equation notation;
    \item recognize graph-theory notation;
    \item identify common mathematical arrows;
    \item distinguish similar-looking symbols;
    \item use notation consistently and unambiguously.
\end{itemize}
}

%%%%%%%%%%%%%%%%%%%%%%%%%%%%%%%%%%%%%%%%%%%%%%%%%%%%%%%%%%%%
\subsection{Purpose of Mathematical Notation}
\label{subsec:purpose-mathematical-notation}
%%%%%%%%%%%%%%%%%%%%%%%%%%%%%%%%%%%%%%%%%%%%%%%%%%%%%%%%%%%%

Mathematical notation compresses complex statements into compact symbolic
form.

For example,

\[
\boxed{
\forall x\in\R,
\qquad
x^2\geq0
}
\]

means:

\[
\boxed{
\text{For every real number }x,\text{ the square }x^2\text{ is nonnegative.}
}
\]

Good notation should be:

\begin{itemize}
    \item consistent;
    \item defined before use;
    \item visually distinguishable;
    \item appropriate to the mathematical context.
\end{itemize}

%%%%%%%%%%%%%%%%%%%%%%%%%%%%%%%%%%%%%%%%%%%%%%%%%%%%%%%%%%%%
\subsection{Basic Arithmetic Symbols}
\label{subsec:basic-arithmetic-symbols}
%%%%%%%%%%%%%%%%%%%%%%%%%%%%%%%%%%%%%%%%%%%%%%%%%%%%%%%%%%%%

\begin{center}
\begin{tabular}{c|l}
\textbf{Symbol} & \textbf{Meaning} \\ \hline
\( + \) & addition \\
\( - \) & subtraction or negation \\
\( \pm \) & plus or minus \\
\( \mp \) & minus or plus \\
\( \times \) & multiplication \\
\( \cdot \) & multiplication or dot product \\
\( \div \) & division \\
\( / \) & division or quotient notation \\
\( = \) & equality \\
\( \neq \) & not equal \\
\( \approx \) & approximately equal \\
\( \equiv \) & identically equal or congruent, depending on context \\
\( \propto \) & proportional to \\
\( \% \) & percent
\end{tabular}
\end{center}

%%%%%%%%%%%%%%%%%%%%%%%%%%%%%%%%%%%%%%%%%%%%%%%%%%%%%%%%%%%%
\subsection{Equality Symbols}
\label{subsec:equality-symbols}
%%%%%%%%%%%%%%%%%%%%%%%%%%%%%%%%%%%%%%%%%%%%%%%%%%%%%%%%%%%%

The symbol

\[
\boxed{
a=b
}
\]

means that \(a\) and \(b\) represent the same mathematical quantity.

The notation

\[
\boxed{
a\neq b
}
\]

means they are not equal.

Approximate equality is written

\[
\boxed{
a\approx b.
}
\]

%%%%%%%%%%%%%%%%%%%%%%%%%%%%%%%%%%%%%%%%%%%%%%%%%%%%%%%%%%%%
\subsection{Definition Symbols}
\label{subsec:definition-symbols}
%%%%%%%%%%%%%%%%%%%%%%%%%%%%%%%%%%%%%%%%%%%%%%%%%%%%%%%%%%%%

Definitions are sometimes written using

\[
\boxed{
:=
}
\]

or

\[
\boxed{
\coloneqq.
}
\]

For example,

\[
\boxed{
f(x)
:=
x^2+1
}
\]

means that the function \(f\) is being defined by this expression.

%%%%%%%%%%%%%%%%%%%%%%%%%%%%%%%%%%%%%%%%%%%%%%%%%%%%%%%%%%%%
\subsection{Inequality Symbols}
\label{subsec:inequality-symbols}
%%%%%%%%%%%%%%%%%%%%%%%%%%%%%%%%%%%%%%%%%%%%%%%%%%%%%%%%%%%%

\begin{center}
\begin{tabular}{c|l}
\textbf{Symbol} & \textbf{Meaning} \\ \hline
\( < \) & less than \\
\( > \) & greater than \\
\( \leq \) & less than or equal to \\
\( \geq \) & greater than or equal to \\
\( \ll \) & much less than \\
\( \gg \) & much greater than
\end{tabular}
\end{center}

For example,

\[
\boxed{
2<5
}
\]

and

\[
\boxed{
x\geq0.
}
\]

%%%%%%%%%%%%%%%%%%%%%%%%%%%%%%%%%%%%%%%%%%%%%%%%%%%%%%%%%%%%
\subsection{Absolute Value}
\label{subsec:absolute-value-notation}
%%%%%%%%%%%%%%%%%%%%%%%%%%%%%%%%%%%%%%%%%%%%%%%%%%%%%%%%%%%%

For a real number \(x\),

\[
\boxed{
|x|
}
\]

denotes absolute value.

\[
|x|
=
\begin{cases}
x, & x\geq0,\\
-x, & x<0.
\end{cases}
\]

Absolute value represents distance from zero.

%%%%%%%%%%%%%%%%%%%%%%%%%%%%%%%%%%%%%%%%%%%%%%%%%%%%%%%%%%%%
\subsection{Norms}
\label{subsec:norm-notation-reference}
%%%%%%%%%%%%%%%%%%%%%%%%%%%%%%%%%%%%%%%%%%%%%%%%%%%%%%%%%%%%

A norm is written

\[
\boxed{
\|\mathbf{x}\|.
}
\]

For Euclidean vectors,

\[
\boxed{
\|\mathbf{x}\|_2
=
\sqrt{
x_1^2+\cdots+x_n^2
}.
}
\]

Other common norms include

\[
\boxed{
\|\mathbf{x}\|_1
=
\sum_{i=1}^{n}|x_i|
}
\]

and

\[
\boxed{
\|\mathbf{x}\|_\infty
=
\max_i|x_i|.
}
\]

%%%%%%%%%%%%%%%%%%%%%%%%%%%%%%%%%%%%%%%%%%%%%%%%%%%%%%%%%%%%
\subsection{Number Sets}
\label{subsec:number-set-symbols}
%%%%%%%%%%%%%%%%%%%%%%%%%%%%%%%%%%%%%%%%%%%%%%%%%%%%%%%%%%%%

\begin{center}
\begin{tabular}{c|l}
\textbf{Symbol} & \textbf{Set} \\ \hline
\( \mathbb{N} \) & natural numbers \\
\( \mathbb{Z} \) & integers \\
\( \mathbb{Q} \) & rational numbers \\
\( \mathbb{R} \) & real numbers \\
\( \mathbb{C} \) & complex numbers
\end{tabular}
\end{center}

Typical containments are

\[
\boxed{
\mathbb{N}
\subset
\mathbb{Z}
\subset
\mathbb{Q}
\subset
\mathbb{R}
\subset
\mathbb{C}.
}
\]

%%%%%%%%%%%%%%%%%%%%%%%%%%%%%%%%%%%%%%%%%%%%%%%%%%%%%%%%%%%%
\subsection{Set Membership}
\label{subsec:set-membership-symbols}
%%%%%%%%%%%%%%%%%%%%%%%%%%%%%%%%%%%%%%%%%%%%%%%%%%%%%%%%%%%%

The notation

\[
\boxed{
x\in A
}
\]

means \(x\) is an element of \(A\).

The notation

\[
\boxed{
x\notin A
}
\]

means \(x\) is not an element of \(A\).

%%%%%%%%%%%%%%%%%%%%%%%%%%%%%%%%%%%%%%%%%%%%%%%%%%%%%%%%%%%%
\subsection{Subsets}
\label{subsec:subset-symbols}
%%%%%%%%%%%%%%%%%%%%%%%%%%%%%%%%%%%%%%%%%%%%%%%%%%%%%%%%%%%%

\[
\boxed{
A\subseteq B
}
\]

means every element of \(A\) belongs to \(B\).

A proper subset may be written

\[
\boxed{
A\subsetneq B.
}
\]

Notation conventions for \(\subset\) vary between texts, so the meaning
should be stated when ambiguity matters.

%%%%%%%%%%%%%%%%%%%%%%%%%%%%%%%%%%%%%%%%%%%%%%%%%%%%%%%%%%%%
\subsection{Set Operations}
\label{subsec:set-operation-symbols}
%%%%%%%%%%%%%%%%%%%%%%%%%%%%%%%%%%%%%%%%%%%%%%%%%%%%%%%%%%%%

\begin{center}
\begin{tabular}{c|l}
\textbf{Symbol} & \textbf{Meaning} \\ \hline
\( A\cup B \) & union \\
\( A\cap B \) & intersection \\
\( A\setminus B \) & set difference \\
\( A^c \) & complement \\
\( \varnothing \) & empty set \\
\( |A| \) & cardinality of a finite set \\
\( \mathcal{P}(A) \) & power set of \(A\)
\end{tabular}
\end{center}

%%%%%%%%%%%%%%%%%%%%%%%%%%%%%%%%%%%%%%%%%%%%%%%%%%%%%%%%%%%%
\subsection{Set Builder Notation}
\label{subsec:set-builder-notation}
%%%%%%%%%%%%%%%%%%%%%%%%%%%%%%%%%%%%%%%%%%%%%%%%%%%%%%%%%%%%

A set may be specified by a property:

\[
\boxed{
\{x\in\R:x>0\}.
}
\]

This is read:

\[
\boxed{
\text{the set of real numbers }x\text{ such that }x>0.
}
\]

The vertical bar may also be used:

\[
\boxed{
\{x\in\R\mid x>0\}.
}
\]

%%%%%%%%%%%%%%%%%%%%%%%%%%%%%%%%%%%%%%%%%%%%%%%%%%%%%%%%%%%%
\subsection{Cartesian Products}
\label{subsec:cartesian-product-notation}
%%%%%%%%%%%%%%%%%%%%%%%%%%%%%%%%%%%%%%%%%%%%%%%%%%%%%%%%%%%%

The Cartesian product of sets \(A\) and \(B\) is

\[
\boxed{
A\times B
=
\{
(a,b):a\in A,\ b\in B
\}.
}
\]

For example,

\[
\boxed{
\R^2
=
\R\times\R.
}
\]

%%%%%%%%%%%%%%%%%%%%%%%%%%%%%%%%%%%%%%%%%%%%%%%%%%%%%%%%%%%%
\subsection{Intervals}
\label{subsec:interval-notation-reference}
%%%%%%%%%%%%%%%%%%%%%%%%%%%%%%%%%%%%%%%%%%%%%%%%%%%%%%%%%%%%

\begin{center}
\begin{tabular}{c|l}
\textbf{Interval} & \textbf{Meaning} \\ \hline
\( (a,b) \) & \(a<x<b\) \\
\( [a,b] \) & \(a\leq x\leq b\) \\
\( [a,b) \) & \(a\leq x<b\) \\
\( (a,b] \) & \(a<x\leq b\) \\
\( (-\infty,a) \) & \(x<a\) \\
\( (a,\infty) \) & \(x>a\)
\end{tabular}
\end{center}

Infinity is never included as an endpoint.

Thus one writes

\[
\boxed{
(-\infty,a]
}
\]

rather than using a square bracket next to \(-\infty\).

%%%%%%%%%%%%%%%%%%%%%%%%%%%%%%%%%%%%%%%%%%%%%%%%%%%%%%%%%%%%
\subsection{Logical Symbols}
\label{subsec:logical-symbols-reference}
%%%%%%%%%%%%%%%%%%%%%%%%%%%%%%%%%%%%%%%%%%%%%%%%%%%%%%%%%%%%

\begin{center}
\begin{tabular}{c|l}
\textbf{Symbol} & \textbf{Name / Meaning} \\ \hline
\( \neg P \) & negation: not \(P\) \\
\( P\land Q \) & conjunction: \(P\) and \(Q\) \\
\( P\lor Q \) & disjunction: \(P\) or \(Q\) \\
\( P\Rightarrow Q \) & implication: if \(P\), then \(Q\) \\
\( P\Leftarrow Q \) & \(P\) if \(Q\) \\
\( P\Leftrightarrow Q \) & if and only if \\
\( \forall \) & for all \\
\( \exists \) & there exists \\
\( \exists! \) & there exists exactly one \\
\( \therefore \) & therefore \\
\( \because \) & because
\end{tabular}
\end{center}

%%%%%%%%%%%%%%%%%%%%%%%%%%%%%%%%%%%%%%%%%%%%%%%%%%%%%%%%%%%%
\subsection{Implication}
\label{subsec:implication-symbol}
%%%%%%%%%%%%%%%%%%%%%%%%%%%%%%%%%%%%%%%%%%%%%%%%%%%%%%%%%%%%

The statement

\[
\boxed{
P\Rightarrow Q
}
\]

means:

\[
\boxed{
\text{if }P\text{ is true, then }Q\text{ is true}.
}
\]

Its converse is

\[
\boxed{
Q\Rightarrow P.
}
\]

These statements are not generally equivalent.

%%%%%%%%%%%%%%%%%%%%%%%%%%%%%%%%%%%%%%%%%%%%%%%%%%%%%%%%%%%%
\subsection{Equivalence}
\label{subsec:logical-equivalence-reference}
%%%%%%%%%%%%%%%%%%%%%%%%%%%%%%%%%%%%%%%%%%%%%%%%%%%%%%%%%%%%

The notation

\[
\boxed{
P\Leftrightarrow Q
}
\]

means that both

\[
P\Rightarrow Q
\]

and

\[
Q\Rightarrow P
\]

hold.

It is read:

\[
\boxed{
P\text{ if and only if }Q.
}
\]

%%%%%%%%%%%%%%%%%%%%%%%%%%%%%%%%%%%%%%%%%%%%%%%%%%%%%%%%%%%%
\subsection{Quantifiers}
\label{subsec:quantifier-symbols}
%%%%%%%%%%%%%%%%%%%%%%%%%%%%%%%%%%%%%%%%%%%%%%%%%%%%%%%%%%%%

The universal quantifier is

\[
\boxed{
\forall.
}
\]

For example,

\[
\boxed{
\forall x\in\R,
\quad
x^2\geq0.
}
\]

The existential quantifier is

\[
\boxed{
\exists.
}
\]

For example,

\[
\boxed{
\exists x\in\R
\text{ such that }
x^2=4.
}
\]

%%%%%%%%%%%%%%%%%%%%%%%%%%%%%%%%%%%%%%%%%%%%%%%%%%%%%%%%%%%%
\subsection{Function Notation}
\label{subsec:function-notation-reference}
%%%%%%%%%%%%%%%%%%%%%%%%%%%%%%%%%%%%%%%%%%%%%%%%%%%%%%%%%%%%

A function may be written

\[
\boxed{
f:A\to B.
}
\]

This means \(f\) maps elements of domain \(A\) to elements of codomain \(B\).

The image of \(x\) is

\[
\boxed{
f(x).
}
\]

%%%%%%%%%%%%%%%%%%%%%%%%%%%%%%%%%%%%%%%%%%%%%%%%%%%%%%%%%%%%
\subsection{Maps-To Arrow}
\label{subsec:maps-to-arrow}
%%%%%%%%%%%%%%%%%%%%%%%%%%%%%%%%%%%%%%%%%%%%%%%%%%%%%%%%%%%%

The notation

\[
\boxed{
x\mapsto f(x)
}
\]

describes the rule of a function.

For example,

\[
\boxed{
f:\R\to\R,
\qquad
x\mapsto x^2.
}
\]

%%%%%%%%%%%%%%%%%%%%%%%%%%%%%%%%%%%%%%%%%%%%%%%%%%%%%%%%%%%%
\subsection{Composition}
\label{subsec:function-composition-reference}
%%%%%%%%%%%%%%%%%%%%%%%%%%%%%%%%%%%%%%%%%%%%%%%%%%%%%%%%%%%%

Function composition is

\[
\boxed{
(f\circ g)(x)
=
f(g(x)).
}
\]

The symbol

\[
\circ
\]

is read ``composed with.''

%%%%%%%%%%%%%%%%%%%%%%%%%%%%%%%%%%%%%%%%%%%%%%%%%%%%%%%%%%%%
\subsection{Inverse Functions}
\label{subsec:inverse-function-notation-reference}
%%%%%%%%%%%%%%%%%%%%%%%%%%%%%%%%%%%%%%%%%%%%%%%%%%%%%%%%%%%%

The inverse of a bijective function \(f\) is written

\[
\boxed{
f^{-1}.
}
\]

It satisfies

\[
\boxed{
f^{-1}(f(x))
=
x.
}
\]

Do not confuse \(f^{-1}(x)\) with

\[
\frac{1}{f(x)}.
\]

%%%%%%%%%%%%%%%%%%%%%%%%%%%%%%%%%%%%%%%%%%%%%%%%%%%%%%%%%%%%
\subsection{Sequences}
\label{subsec:sequence-notation-reference}
%%%%%%%%%%%%%%%%%%%%%%%%%%%%%%%%%%%%%%%%%%%%%%%%%%%%%%%%%%%%

A sequence may be written

\[
\boxed{
(a_n)_{n=1}^{\infty}
}
\]

or simply

\[
\boxed{
(a_n).
}
\]

Its \(n\)-th term is

\[
\boxed{
a_n.
}
\]

%%%%%%%%%%%%%%%%%%%%%%%%%%%%%%%%%%%%%%%%%%%%%%%%%%%%%%%%%%%%
\subsection{Summation Notation}
\label{subsec:summation-notation-reference}
%%%%%%%%%%%%%%%%%%%%%%%%%%%%%%%%%%%%%%%%%%%%%%%%%%%%%%%%%%%%

The Greek capital letter sigma,

\[
\boxed{
\sum
}
\]

denotes summation.

For example,

\[
\boxed{
\sum_{k=1}^{n}a_k
=
a_1+a_2+\cdots+a_n.
}
\]

%%%%%%%%%%%%%%%%%%%%%%%%%%%%%%%%%%%%%%%%%%%%%%%%%%%%%%%%%%%%
\subsection{Product Notation}
\label{subsec:product-notation-reference}
%%%%%%%%%%%%%%%%%%%%%%%%%%%%%%%%%%%%%%%%%%%%%%%%%%%%%%%%%%%%

The Greek capital letter pi,

\[
\boxed{
\prod
}
\]

denotes products.

\[
\boxed{
\prod_{k=1}^{n}a_k
=
a_1a_2\cdots a_n.
}
\]

%%%%%%%%%%%%%%%%%%%%%%%%%%%%%%%%%%%%%%%%%%%%%%%%%%%%%%%%%%%%
\subsection{Factorial}
\label{subsec:factorial-notation-reference}
%%%%%%%%%%%%%%%%%%%%%%%%%%%%%%%%%%%%%%%%%%%%%%%%%%%%%%%%%%%%

For \(n\in\mathbb{N}\),

\[
\boxed{
n!
=
n(n-1)(n-2)\cdots2\cdot1.
}
\]

By convention,

\[
\boxed{
0!=1.
}
\]

%%%%%%%%%%%%%%%%%%%%%%%%%%%%%%%%%%%%%%%%%%%%%%%%%%%%%%%%%%%%
\subsection{Binomial Coefficients}
\label{subsec:binomial-coefficient-notation-reference}
%%%%%%%%%%%%%%%%%%%%%%%%%%%%%%%%%%%%%%%%%%%%%%%%%%%%%%%%%%%%

The number of ways to choose \(k\) objects from \(n\) is

\[
\boxed{
\binom{n}{k}
=
\frac{n!}{k!(n-k)!}.
}
\]

It may also be read as

\[
\boxed{
n\text{ choose }k.
}
\]

%%%%%%%%%%%%%%%%%%%%%%%%%%%%%%%%%%%%%%%%%%%%%%%%%%%%%%%%%%%%
\subsection{Limits}
\label{subsec:limit-notation-reference}
%%%%%%%%%%%%%%%%%%%%%%%%%%%%%%%%%%%%%%%%%%%%%%%%%%%%%%%%%%%%

The notation

\[
\boxed{
\lim_{x\to a}f(x)
=
L
}
\]

means:

\[
\boxed{
\text{the limit as }x\text{ approaches }a\text{ of }f(x)\text{ is }L.
}
\]

One-sided limits are

\[
\boxed{
\lim_{x\to a^-}f(x)
}
\]

and

\[
\boxed{
\lim_{x\to a^+}f(x).
}
\]

%%%%%%%%%%%%%%%%%%%%%%%%%%%%%%%%%%%%%%%%%%%%%%%%%%%%%%%%%%%%
\subsection{Infinity}
\label{subsec:infinity-symbol-reference}
%%%%%%%%%%%%%%%%%%%%%%%%%%%%%%%%%%%%%%%%%%%%%%%%%%%%%%%%%%%%

The symbol

\[
\boxed{
\infty
}
\]

represents unbounded growth or an infinite extent.

It is not an ordinary real number.

For example,

\[
\boxed{
\lim_{x\to\infty}\frac{1}{x}=0.
}
\]

%%%%%%%%%%%%%%%%%%%%%%%%%%%%%%%%%%%%%%%%%%%%%%%%%%%%%%%%%%%%
\subsection{Derivative Notation}
\label{subsec:derivative-notation-reference}
%%%%%%%%%%%%%%%%%%%%%%%%%%%%%%%%%%%%%%%%%%%%%%%%%%%%%%%%%%%%

Common derivative notation includes

\[
\boxed{
f'(x)
}
\]

\[
\boxed{
\frac{df}{dx}
}
\]

and

\[
\boxed{
D_xf.
}
\]

For a function \(y=f(x)\),

\[
\boxed{
\frac{dy}{dx}
=
f'(x).
}
\]

%%%%%%%%%%%%%%%%%%%%%%%%%%%%%%%%%%%%%%%%%%%%%%%%%%%%%%%%%%%%
\subsection{Higher Derivatives}
\label{subsec:higher-derivative-notation-reference}
%%%%%%%%%%%%%%%%%%%%%%%%%%%%%%%%%%%%%%%%%%%%%%%%%%%%%%%%%%%%

The second derivative may be written

\[
\boxed{
f''(x)
}
\]

or

\[
\boxed{
\frac{d^2f}{dx^2}.
}
\]

The \(n\)-th derivative is

\[
\boxed{
f^{(n)}(x).
}
\]

%%%%%%%%%%%%%%%%%%%%%%%%%%%%%%%%%%%%%%%%%%%%%%%%%%%%%%%%%%%%
\subsection{Partial Derivatives}
\label{subsec:partial-derivative-notation-reference}
%%%%%%%%%%%%%%%%%%%%%%%%%%%%%%%%%%%%%%%%%%%%%%%%%%%%%%%%%%%%

For

\[
f(x,y),
\]

the partial derivative with respect to \(x\) is

\[
\boxed{
\frac{\partial f}{\partial x}.
}
\]

A second mixed partial is

\[
\boxed{
\frac{\partial^2f}
{\partial y\partial x}.
}
\]

%%%%%%%%%%%%%%%%%%%%%%%%%%%%%%%%%%%%%%%%%%%%%%%%%%%%%%%%%%%%
\subsection{Gradient}
\label{subsec:gradient-notation-reference}
%%%%%%%%%%%%%%%%%%%%%%%%%%%%%%%%%%%%%%%%%%%%%%%%%%%%%%%%%%%%

For scalar field

\[
f:\R^n\to\R,
\]

the gradient is

\[
\boxed{
\nabla f
=
\begin{bmatrix}
\frac{\partial f}{\partial x_1}\\
\vdots\\
\frac{\partial f}{\partial x_n}
\end{bmatrix}.
}
\]

The symbol

\[
\nabla
\]

is called nabla or del.

%%%%%%%%%%%%%%%%%%%%%%%%%%%%%%%%%%%%%%%%%%%%%%%%%%%%%%%%%%%%
\subsection{Divergence}
\label{subsec:divergence-notation-reference}
%%%%%%%%%%%%%%%%%%%%%%%%%%%%%%%%%%%%%%%%%%%%%%%%%%%%%%%%%%%%

For vector field

\[
\mathbf{F}
=
(F_1,\ldots,F_n),
\]

the divergence is

\[
\boxed{
\nabla\cdot\mathbf{F}
=
\sum_{i=1}^{n}
\frac{\partial F_i}{\partial x_i}.
}
\]

%%%%%%%%%%%%%%%%%%%%%%%%%%%%%%%%%%%%%%%%%%%%%%%%%%%%%%%%%%%%
\subsection{Curl}
\label{subsec:curl-notation-reference}
%%%%%%%%%%%%%%%%%%%%%%%%%%%%%%%%%%%%%%%%%%%%%%%%%%%%%%%%%%%%

In three dimensions,

\[
\boxed{
\nabla\times\mathbf{F}
}
\]

denotes curl.

For

\[
\mathbf{F}
=
(F_1,F_2,F_3),
\]

\[
\nabla\times\mathbf{F}
=
\begin{bmatrix}
F_{3,y}-F_{2,z}\\
F_{1,z}-F_{3,x}\\
F_{2,x}-F_{1,y}
\end{bmatrix}.
\]

%%%%%%%%%%%%%%%%%%%%%%%%%%%%%%%%%%%%%%%%%%%%%%%%%%%%%%%%%%%%
\subsection{Laplacian}
\label{subsec:laplacian-notation-reference}
%%%%%%%%%%%%%%%%%%%%%%%%%%%%%%%%%%%%%%%%%%%%%%%%%%%%%%%%%%%%

The Laplacian of a scalar function is

\[
\boxed{
\Delta f
=
\nabla^2f.
}
\]

In Cartesian coordinates,

\[
\boxed{
\Delta f
=
\sum_{i=1}^{n}
\frac{\partial^2f}{\partial x_i^2}.
}
\]

%%%%%%%%%%%%%%%%%%%%%%%%%%%%%%%%%%%%%%%%%%%%%%%%%%%%%%%%%%%%
\subsection{Integrals}
\label{subsec:integral-notation-reference}
%%%%%%%%%%%%%%%%%%%%%%%%%%%%%%%%%%%%%%%%%%%%%%%%%%%%%%%%%%%%

An indefinite integral is

\[
\boxed{
\int f(x)\,dx.
}
\]

A definite integral is

\[
\boxed{
\int_a^bf(x)\,dx.
}
\]

Multiple integrals include

\[
\boxed{
\iint_Df(x,y)\,dA
}
\]

and

\[
\boxed{
\iiint_Vf(x,y,z)\,dV.
}
\]

%%%%%%%%%%%%%%%%%%%%%%%%%%%%%%%%%%%%%%%%%%%%%%%%%%%%%%%%%%%%
\subsection{Line and Surface Integrals}
\label{subsec:line-surface-integral-notation}
%%%%%%%%%%%%%%%%%%%%%%%%%%%%%%%%%%%%%%%%%%%%%%%%%%%%%%%%%%%%

A line integral may be written

\[
\boxed{
\int_C
\mathbf{F}\cdot d\mathbf{r}.
}
\]

A surface flux integral may be written

\[
\boxed{
\iint_S
\mathbf{F}\cdot\mathbf{n}\,dS.
}
\]

A closed contour integral is often written

\[
\boxed{
\oint_C.
}
\]

%%%%%%%%%%%%%%%%%%%%%%%%%%%%%%%%%%%%%%%%%%%%%%%%%%%%%%%%%%%%
\subsection{Differentials}
\label{subsec:differential-symbol-reference}
%%%%%%%%%%%%%%%%%%%%%%%%%%%%%%%%%%%%%%%%%%%%%%%%%%%%%%%%%%%%

The symbol

\[
\boxed{
dx
}
\]

appears in derivative and integral notation.

Examples include

\[
\boxed{
dy
=
f'(x)\,dx
}
\]

and

\[
\boxed{
\int f(x)\,dx.
}
\]

Its formal interpretation depends on context.

%%%%%%%%%%%%%%%%%%%%%%%%%%%%%%%%%%%%%%%%%%%%%%%%%%%%%%%%%%%%
\subsection{Vectors}
\label{subsec:vector-notation-reference}
%%%%%%%%%%%%%%%%%%%%%%%%%%%%%%%%%%%%%%%%%%%%%%%%%%%%%%%%%%%%

A vector may be written

\[
\boxed{
\mathbf{v}
}
\]

or sometimes

\[
\boxed{
\vec{v}.
}
\]

In coordinates,

\[
\boxed{
\mathbf{v}
=
\begin{bmatrix}
v_1\\
\vdots\\
v_n
\end{bmatrix}.
}
\]

%%%%%%%%%%%%%%%%%%%%%%%%%%%%%%%%%%%%%%%%%%%%%%%%%%%%%%%%%%%%
\subsection{Dot Product}
\label{subsec:dot-product-notation-reference}
%%%%%%%%%%%%%%%%%%%%%%%%%%%%%%%%%%%%%%%%%%%%%%%%%%%%%%%%%%%%

The dot product is

\[
\boxed{
\mathbf{u}\cdot\mathbf{v}
}
\]

or

\[
\boxed{
\mathbf{u}^T\mathbf{v}.
}
\]

In \(\R^n\),

\[
\boxed{
\mathbf{u}\cdot\mathbf{v}
=
\sum_{i=1}^{n}u_iv_i.
}
\]

%%%%%%%%%%%%%%%%%%%%%%%%%%%%%%%%%%%%%%%%%%%%%%%%%%%%%%%%%%%%
\subsection{Cross Product}
\label{subsec:cross-product-reference}
%%%%%%%%%%%%%%%%%%%%%%%%%%%%%%%%%%%%%%%%%%%%%%%%%%%%%%%%%%%%

In \(\R^3\),

\[
\boxed{
\mathbf{u}\times\mathbf{v}
}
\]

denotes the cross product.

The result is perpendicular to both vectors.

%%%%%%%%%%%%%%%%%%%%%%%%%%%%%%%%%%%%%%%%%%%%%%%%%%%%%%%%%%%%
\subsection{Matrices}
\label{subsec:matrix-notation-reference}
%%%%%%%%%%%%%%%%%%%%%%%%%%%%%%%%%%%%%%%%%%%%%%%%%%%%%%%%%%%%

A matrix is commonly denoted by a capital letter such as

\[
\boxed{
A.
}
\]

Its entry in row \(i\), column \(j\) is

\[
\boxed{
A_{ij}.
}
\]

An \(m\times n\) matrix belongs to

\[
\boxed{
\R^{m\times n}.
}
\]

%%%%%%%%%%%%%%%%%%%%%%%%%%%%%%%%%%%%%%%%%%%%%%%%%%%%%%%%%%%%
\subsection{Transpose}
\label{subsec:transpose-symbol-reference}
%%%%%%%%%%%%%%%%%%%%%%%%%%%%%%%%%%%%%%%%%%%%%%%%%%%%%%%%%%%%

The transpose of \(A\) is

\[
\boxed{
A^T.
}
\]

Its entries satisfy

\[
\boxed{
(A^T)_{ij}
=
A_{ji}.
}
\]

%%%%%%%%%%%%%%%%%%%%%%%%%%%%%%%%%%%%%%%%%%%%%%%%%%%%%%%%%%%%
\subsection{Inverse Matrix}
\label{subsec:inverse-matrix-reference}
%%%%%%%%%%%%%%%%%%%%%%%%%%%%%%%%%%%%%%%%%%%%%%%%%%%%%%%%%%%%

If \(A\) is invertible,

\[
\boxed{
A^{-1}
}
\]

satisfies

\[
\boxed{
AA^{-1}
=
A^{-1}A
=
I.
}
\]

The symbol \(I\) denotes the identity matrix.

%%%%%%%%%%%%%%%%%%%%%%%%%%%%%%%%%%%%%%%%%%%%%%%%%%%%%%%%%%%%
\subsection{Determinant}
\label{subsec:determinant-symbol-reference}
%%%%%%%%%%%%%%%%%%%%%%%%%%%%%%%%%%%%%%%%%%%%%%%%%%%%%%%%%%%%

The determinant is written

\[
\boxed{
\det(A)
}
\]

or

\[
\boxed{
|A|.
}
\]

For a \(2\times2\) matrix,

\[
\boxed{
\det
\begin{bmatrix}
a & b\\
c & d
\end{bmatrix}
=
ad-bc.
}
\]

%%%%%%%%%%%%%%%%%%%%%%%%%%%%%%%%%%%%%%%%%%%%%%%%%%%%%%%%%%%%
\subsection{Trace}
\label{subsec:trace-symbol-reference}
%%%%%%%%%%%%%%%%%%%%%%%%%%%%%%%%%%%%%%%%%%%%%%%%%%%%%%%%%%%%

The trace is

\[
\boxed{
\operatorname{tr}(A)
=
\sum_iA_{ii}.
}
\]

It is the sum of diagonal entries.

%%%%%%%%%%%%%%%%%%%%%%%%%%%%%%%%%%%%%%%%%%%%%%%%%%%%%%%%%%%%
\subsection{Rank}
\label{subsec:rank-symbol-reference}
%%%%%%%%%%%%%%%%%%%%%%%%%%%%%%%%%%%%%%%%%%%%%%%%%%%%%%%%%%%%

The rank of a matrix is written

\[
\boxed{
\operatorname{rank}(A).
}
\]

It equals the dimension of the column space and row space.

%%%%%%%%%%%%%%%%%%%%%%%%%%%%%%%%%%%%%%%%%%%%%%%%%%%%%%%%%%%%
\subsection{Eigenvalues and Eigenvectors}
\label{subsec:eigenvalue-notation-reference}
%%%%%%%%%%%%%%%%%%%%%%%%%%%%%%%%%%%%%%%%%%%%%%%%%%%%%%%%%%%%

An eigenvalue-eigenvector pair satisfies

\[
\boxed{
A\mathbf{v}
=
\lambda\mathbf{v},
}
\]

where

\[
\mathbf{v}\neq0.
\]

The scalar

\[
\lambda
\]

is an eigenvalue.

The vector

\[
\mathbf{v}
\]

is an eigenvector.

%%%%%%%%%%%%%%%%%%%%%%%%%%%%%%%%%%%%%%%%%%%%%%%%%%%%%%%%%%%%
\subsection{Identity Matrix}
\label{subsec:identity-matrix-reference}
%%%%%%%%%%%%%%%%%%%%%%%%%%%%%%%%%%%%%%%%%%%%%%%%%%%%%%%%%%%%

The identity matrix is

\[
\boxed{
I
=
\begin{bmatrix}
1 & 0 & \cdots & 0\\
0 & 1 & \cdots & 0\\
\vdots & \vdots & \ddots & \vdots\\
0 & 0 & \cdots & 1
\end{bmatrix}.
}
\]

It satisfies

\[
\boxed{
AI=IA=A.
}
\]

%%%%%%%%%%%%%%%%%%%%%%%%%%%%%%%%%%%%%%%%%%%%%%%%%%%%%%%%%%%%
\subsection{Zero Vector and Zero Matrix}
\label{subsec:zero-vector-matrix-reference}
%%%%%%%%%%%%%%%%%%%%%%%%%%%%%%%%%%%%%%%%%%%%%%%%%%%%%%%%%%%%

The symbol

\[
\boxed{
\mathbf{0}
}
\]

usually denotes a zero vector.

The symbol

\[
\boxed{
0
}
\]

may denote a scalar zero or zero matrix depending on context.

Dimensions should be clear from the equation.

%%%%%%%%%%%%%%%%%%%%%%%%%%%%%%%%%%%%%%%%%%%%%%%%%%%%%%%%%%%%
\subsection{Complex Numbers}
\label{subsec:complex-symbols-reference}
%%%%%%%%%%%%%%%%%%%%%%%%%%%%%%%%%%%%%%%%%%%%%%%%%%%%%%%%%%%%

A complex number is

\[
\boxed{
z
=
a+bi,
}
\]

where

\[
\boxed{
i^2=-1.
}
\]

The real and imaginary parts are

\[
\boxed{
\operatorname{Re}(z)=a,
}
\]

\[
\boxed{
\operatorname{Im}(z)=b.
}
\]

%%%%%%%%%%%%%%%%%%%%%%%%%%%%%%%%%%%%%%%%%%%%%%%%%%%%%%%%%%%%
\subsection{Complex Conjugate}
\label{subsec:complex-conjugate-reference}
%%%%%%%%%%%%%%%%%%%%%%%%%%%%%%%%%%%%%%%%%%%%%%%%%%%%%%%%%%%%

The complex conjugate of

\[
z=a+bi
\]

is

\[
\boxed{
\overline{z}
=
a-bi.
}
\]

Its modulus satisfies

\[
\boxed{
|z|^2
=
z\overline{z}.
}
\]

%%%%%%%%%%%%%%%%%%%%%%%%%%%%%%%%%%%%%%%%%%%%%%%%%%%%%%%%%%%%
\subsection{Argument of a Complex Number}
\label{subsec:complex-argument-reference}
%%%%%%%%%%%%%%%%%%%%%%%%%%%%%%%%%%%%%%%%%%%%%%%%%%%%%%%%%%%%

For

\[
z=re^{i\theta},
\]

the angle \(\theta\) is an argument of \(z\).

It is written

\[
\boxed{
\arg(z).
}
\]

A principal argument is often written

\[
\boxed{
\operatorname{Arg}(z).
}
\]

%%%%%%%%%%%%%%%%%%%%%%%%%%%%%%%%%%%%%%%%%%%%%%%%%%%%%%%%%%%%
\subsection{Probability Symbols}
\label{subsec:probability-symbols-reference}
%%%%%%%%%%%%%%%%%%%%%%%%%%%%%%%%%%%%%%%%%%%%%%%%%%%%%%%%%%%%

\begin{center}
\begin{tabular}{c|l}
\textbf{Symbol} & \textbf{Meaning} \\ \hline
\( \Omega \) & sample space \\
\( A,B \) & events \\
\( \Pr(A) \) & probability of event \(A\) \\
\( A^c \) & complement of \(A\) \\
\( A\cap B \) & both events occur \\
\( A\cup B \) & at least one occurs \\
\( X \) & random variable \\
\( X\sim F \) & \(X\) has distribution \(F\)
\end{tabular}
\end{center}

%%%%%%%%%%%%%%%%%%%%%%%%%%%%%%%%%%%%%%%%%%%%%%%%%%%%%%%%%%%%
\subsection{Conditional Probability}
\label{subsec:conditional-probability-reference}
%%%%%%%%%%%%%%%%%%%%%%%%%%%%%%%%%%%%%%%%%%%%%%%%%%%%%%%%%%%%

Conditional probability is

\[
\boxed{
\Pr(A\mid B)
=
\frac{
\Pr(A\cap B)
}{
\Pr(B)
}
}
\]

provided

\[
\Pr(B)>0.
\]

The vertical bar is read ``given.''

%%%%%%%%%%%%%%%%%%%%%%%%%%%%%%%%%%%%%%%%%%%%%%%%%%%%%%%%%%%%
\subsection{Independence}
\label{subsec:independence-notation-reference}
%%%%%%%%%%%%%%%%%%%%%%%%%%%%%%%%%%%%%%%%%%%%%%%%%%%%%%%%%%%%

Events \(A\) and \(B\) are independent when

\[
\boxed{
\Pr(A\cap B)
=
\Pr(A)\Pr(B).
}
\]

Random-variable independence is expressed similarly through joint
distributions.

%%%%%%%%%%%%%%%%%%%%%%%%%%%%%%%%%%%%%%%%%%%%%%%%%%%%%%%%%%%%
\subsection{Expectation}
\label{subsec:expectation-notation-reference}
%%%%%%%%%%%%%%%%%%%%%%%%%%%%%%%%%%%%%%%%%%%%%%%%%%%%%%%%%%%%

Expected value is written

\[
\boxed{
\mathbb{E}[X].
}
\]

For a discrete random variable,

\[
\boxed{
\mathbb{E}[X]
=
\sum_x
x\Pr(X=x).
}
\]

For a continuous random variable with density \(f_X\),

\[
\boxed{
\mathbb{E}[X]
=
\int_{-\infty}^{\infty}
x f_X(x)\,dx.
}
\]

%%%%%%%%%%%%%%%%%%%%%%%%%%%%%%%%%%%%%%%%%%%%%%%%%%%%%%%%%%%%
\subsection{Variance}
\label{subsec:variance-notation-reference}
%%%%%%%%%%%%%%%%%%%%%%%%%%%%%%%%%%%%%%%%%%%%%%%%%%%%%%%%%%%%

Variance is

\[
\boxed{
\operatorname{Var}(X).
}
\]

Equivalent notation includes

\[
\boxed{
\sigma_X^2.
}
\]

It is defined by

\[
\boxed{
\operatorname{Var}(X)
=
\mathbb{E}
\left[
(X-\mathbb{E}[X])^2
\right].
}
\]

%%%%%%%%%%%%%%%%%%%%%%%%%%%%%%%%%%%%%%%%%%%%%%%%%%%%%%%%%%%%
\subsection{Covariance}
\label{subsec:covariance-notation-reference}
%%%%%%%%%%%%%%%%%%%%%%%%%%%%%%%%%%%%%%%%%%%%%%%%%%%%%%%%%%%%

Covariance is

\[
\boxed{
\operatorname{Cov}(X,Y).
}
\]

It is

\[
\boxed{
\operatorname{Cov}(X,Y)
=
\mathbb{E}
\left[
(X-\mathbb{E}[X])
(Y-\mathbb{E}[Y])
\right].
}
\]

%%%%%%%%%%%%%%%%%%%%%%%%%%%%%%%%%%%%%%%%%%%%%%%%%%%%%%%%%%%%
\subsection{Correlation}
\label{subsec:correlation-notation-reference}
%%%%%%%%%%%%%%%%%%%%%%%%%%%%%%%%%%%%%%%%%%%%%%%%%%%%%%%%%%%%

Correlation is often written

\[
\boxed{
\rho_{X,Y}
}
\]

and defined by

\[
\boxed{
\rho_{X,Y}
=
\frac{
\operatorname{Cov}(X,Y)
}{
\sigma_X\sigma_Y
}.
}
\]

When variances are positive,

\[
\boxed{
-1\leq\rho_{X,Y}\leq1.
}
\]

%%%%%%%%%%%%%%%%%%%%%%%%%%%%%%%%%%%%%%%%%%%%%%%%%%%%%%%%%%%%
\subsection{Probability Density and Distribution Functions}
\label{subsec:density-distribution-reference}
%%%%%%%%%%%%%%%%%%%%%%%%%%%%%%%%%%%%%%%%%%%%%%%%%%%%%%%%%%%%

A probability density is commonly written

\[
\boxed{
f_X(x).
}
\]

A cumulative distribution function is

\[
\boxed{
F_X(x)
=
\Pr(X\leq x).
}
\]

A probability mass function may be written

\[
\boxed{
p_X(x)
=
\Pr(X=x).
}
\]

%%%%%%%%%%%%%%%%%%%%%%%%%%%%%%%%%%%%%%%%%%%%%%%%%%%%%%%%%%%%
\subsection{Statistical Symbols}
\label{subsec:statistical-symbols-reference}
%%%%%%%%%%%%%%%%%%%%%%%%%%%%%%%%%%%%%%%%%%%%%%%%%%%%%%%%%%%%

Common notation includes

\[
\boxed{
\bar{x}
}
\]

for sample mean,

\[
\boxed{
s^2
}
\]

for sample variance,

\[
\boxed{
\hat{\theta}
}
\]

for an estimator of parameter \(\theta\),

and

\[
\boxed{
n
}
\]

for sample size.

%%%%%%%%%%%%%%%%%%%%%%%%%%%%%%%%%%%%%%%%%%%%%%%%%%%%%%%%%%%%
\subsection{Sample Mean}
\label{subsec:sample-mean-reference}
%%%%%%%%%%%%%%%%%%%%%%%%%%%%%%%%%%%%%%%%%%%%%%%%%%%%%%%%%%%%

For observations

\[
x_1,\ldots,x_n,
\]

the sample mean is

\[
\boxed{
\bar{x}
=
\frac1n
\sum_{i=1}^{n}x_i.
}
\]

%%%%%%%%%%%%%%%%%%%%%%%%%%%%%%%%%%%%%%%%%%%%%%%%%%%%%%%%%%%%
\subsection{Sample Variance}
\label{subsec:sample-variance-reference}
%%%%%%%%%%%%%%%%%%%%%%%%%%%%%%%%%%%%%%%%%%%%%%%%%%%%%%%%%%%%

A common unbiased sample variance is

\[
\boxed{
s^2
=
\frac{1}{n-1}
\sum_{i=1}^{n}
(x_i-\bar{x})^2.
}
\]

The sample standard deviation is

\[
\boxed{
s=\sqrt{s^2}.
}
\]

%%%%%%%%%%%%%%%%%%%%%%%%%%%%%%%%%%%%%%%%%%%%%%%%%%%%%%%%%%%%
\subsection{Estimator Notation}
\label{subsec:estimator-notation-reference}
%%%%%%%%%%%%%%%%%%%%%%%%%%%%%%%%%%%%%%%%%%%%%%%%%%%%%%%%%%%%

A hat usually denotes an estimate:

\[
\boxed{
\hat{\theta}.
}
\]

For example,

\[
\boxed{
\hat{\beta}
}
\]

may denote an estimated regression coefficient.

%%%%%%%%%%%%%%%%%%%%%%%%%%%%%%%%%%%%%%%%%%%%%%%%%%%%%%%%%%%%
\subsection{Optimization Symbols}
\label{subsec:optimization-symbols-reference}
%%%%%%%%%%%%%%%%%%%%%%%%%%%%%%%%%%%%%%%%%%%%%%%%%%%%%%%%%%%%

A minimization problem is written

\[
\boxed{
\min_{x\in X}
f(x).
}
\]

A maximization problem is

\[
\boxed{
\max_{x\in X}
f(x).
}
\]

An optimizer is often written

\[
\boxed{
x^*
}
\]

or

\[
\boxed{
x^\star.
}
\]

%%%%%%%%%%%%%%%%%%%%%%%%%%%%%%%%%%%%%%%%%%%%%%%%%%%%%%%%%%%%
\subsection{Argmin and Argmax}
\label{subsec:argmin-argmax-reference}
%%%%%%%%%%%%%%%%%%%%%%%%%%%%%%%%%%%%%%%%%%%%%%%%%%%%%%%%%%%%

The minimum value is

\[
\min_x f(x).
\]

The set of points where the minimum is attained is

\[
\boxed{
\operatorname*{arg\,min}_x
f(x).
}
\]

Similarly,

\[
\boxed{
\operatorname*{arg\,max}_x
f(x).
}
\]

%%%%%%%%%%%%%%%%%%%%%%%%%%%%%%%%%%%%%%%%%%%%%%%%%%%%%%%%%%%%
\subsection{Gradient and Hessian in Optimization}
\label{subsec:optimization-gradient-hessian-reference}
%%%%%%%%%%%%%%%%%%%%%%%%%%%%%%%%%%%%%%%%%%%%%%%%%%%%%%%%%%%%

For scalar objective \(f:\R^n\to\R\),

\[
\boxed{
\nabla f(\mathbf{x})
}
\]

is the gradient.

The Hessian is

\[
\boxed{
\nabla^2f(\mathbf{x})
}
\]

with entries

\[
\boxed{
\left[
\nabla^2f
\right]_{ij}
=
\frac{\partial^2f}
{\partial x_i\partial x_j}.
}
\]

%%%%%%%%%%%%%%%%%%%%%%%%%%%%%%%%%%%%%%%%%%%%%%%%%%%%%%%%%%%%
\subsection{Lagrangian Notation}
\label{subsec:lagrangian-notation-reference}
%%%%%%%%%%%%%%%%%%%%%%%%%%%%%%%%%%%%%%%%%%%%%%%%%%%%%%%%%%%%

For constrained optimization,

\[
\boxed{
\mathcal{L}(x,\lambda)
=
f(x)
+
\lambda g(x)
}
\]

is a simple Lagrangian.

For multiple constraints,

\[
\boxed{
\mathcal{L}
=
f
+
\sum_i\lambda_i g_i
+
\sum_j\nu_j h_j.
}
\]

%%%%%%%%%%%%%%%%%%%%%%%%%%%%%%%%%%%%%%%%%%%%%%%%%%%%%%%%%%%%
\subsection{Differential Equation Notation}
\label{subsec:differential-equation-symbols-reference}
%%%%%%%%%%%%%%%%%%%%%%%%%%%%%%%%%%%%%%%%%%%%%%%%%%%%%%%%%%%%

A first-order ordinary differential equation may be written

\[
\boxed{
\frac{dy}{dt}
=
f(t,y)
}
\]

or

\[
\boxed{
y'=f(t,y).
}
\]

A second-order ODE may be written

\[
\boxed{
y''+ay'+by=0.
}
\]

%%%%%%%%%%%%%%%%%%%%%%%%%%%%%%%%%%%%%%%%%%%%%%%%%%%%%%%%%%%%
\subsection{Dot Notation}
\label{subsec:dot-notation-reference}
%%%%%%%%%%%%%%%%%%%%%%%%%%%%%%%%%%%%%%%%%%%%%%%%%%%%%%%%%%%%

When time is the independent variable,

\[
\boxed{
\dot{x}
=
\frac{dx}{dt}
}
\]

and

\[
\boxed{
\ddot{x}
=
\frac{d^2x}{dt^2}.
}
\]

This notation is common in mechanics and dynamical systems.

%%%%%%%%%%%%%%%%%%%%%%%%%%%%%%%%%%%%%%%%%%%%%%%%%%%%%%%%%%%%
\subsection{Partial Differential Equation Notation}
\label{subsec:pde-notation-reference}
%%%%%%%%%%%%%%%%%%%%%%%%%%%%%%%%%%%%%%%%%%%%%%%%%%%%%%%%%%%%

For

\[
u=u(x,t),
\]

one may write

\[
\boxed{
u_t
=
\frac{\partial u}{\partial t},
}
\]

\[
\boxed{
u_x
=
\frac{\partial u}{\partial x},
}
\]

and

\[
\boxed{
u_{xx}
=
\frac{\partial^2u}{\partial x^2}.
}
\]

%%%%%%%%%%%%%%%%%%%%%%%%%%%%%%%%%%%%%%%%%%%%%%%%%%%%%%%%%%%%
\subsection{Asymptotic Notation}
\label{subsec:asymptotic-notation-reference}
%%%%%%%%%%%%%%%%%%%%%%%%%%%%%%%%%%%%%%%%%%%%%%%%%%%%%%%%%%%%

Asymptotic notation compares growth rates.

Common symbols are:

\[
\boxed{
O(g(n))
}
\]

\[
\boxed{
o(g(n))
}
\]

\[
\boxed{
\Theta(g(n))
}
\]

and

\[
\boxed{
\Omega(g(n)).
}
\]

Their precise meanings depend on the asymptotic context.

%%%%%%%%%%%%%%%%%%%%%%%%%%%%%%%%%%%%%%%%%%%%%%%%%%%%%%%%%%%%
\subsection{Big-O Notation}
\label{subsec:big-o-reference}
%%%%%%%%%%%%%%%%%%%%%%%%%%%%%%%%%%%%%%%%%%%%%%%%%%%%%%%%%%%%

The statement

\[
\boxed{
f(n)=O(g(n))
}
\]

means that for sufficiently large \(n\),

\[
\boxed{
|f(n)|
\leq
C|g(n)|
}
\]

for some constant \(C>0\).

It gives an asymptotic upper bound.

%%%%%%%%%%%%%%%%%%%%%%%%%%%%%%%%%%%%%%%%%%%%%%%%%%%%%%%%%%%%
\subsection{Little-o Notation}
\label{subsec:little-o-reference}
%%%%%%%%%%%%%%%%%%%%%%%%%%%%%%%%%%%%%%%%%%%%%%%%%%%%%%%%%%%%

The statement

\[
\boxed{
f(n)=o(g(n))
}
\]

means

\[
\boxed{
\frac{f(n)}{g(n)}
\to0
}
\]

under the relevant limit.

Thus \(f\) grows strictly more slowly than \(g\) in this asymptotic sense.

%%%%%%%%%%%%%%%%%%%%%%%%%%%%%%%%%%%%%%%%%%%%%%%%%%%%%%%%%%%%
\subsection{Asymptotic Equivalence}
\label{subsec:asymptotic-equivalence-reference}
%%%%%%%%%%%%%%%%%%%%%%%%%%%%%%%%%%%%%%%%%%%%%%%%%%%%%%%%%%%%

The notation

\[
\boxed{
f(x)\sim g(x)
}
\]

often means

\[
\boxed{
\frac{f(x)}{g(x)}
\to1.
}
\]

It should not be confused with informal approximation notation.

%%%%%%%%%%%%%%%%%%%%%%%%%%%%%%%%%%%%%%%%%%%%%%%%%%%%%%%%%%%%
\subsection{Congruence}
\label{subsec:congruence-symbol-reference}
%%%%%%%%%%%%%%%%%%%%%%%%%%%%%%%%%%%%%%%%%%%%%%%%%%%%%%%%%%%%

In number theory,

\[
\boxed{
a\equiv b\pmod n
}
\]

means

\[
\boxed{
n\mid(a-b).
}
\]

The symbol \(\equiv\) can have different meanings elsewhere, so context
matters.

%%%%%%%%%%%%%%%%%%%%%%%%%%%%%%%%%%%%%%%%%%%%%%%%%%%%%%%%%%%%
\subsection{Divisibility}
\label{subsec:divisibility-symbol-reference}
%%%%%%%%%%%%%%%%%%%%%%%%%%%%%%%%%%%%%%%%%%%%%%%%%%%%%%%%%%%%

The notation

\[
\boxed{
a\mid b
}
\]

means \(a\) divides \(b\).

The notation

\[
\boxed{
a\nmid b
}
\]

means \(a\) does not divide \(b\).

%%%%%%%%%%%%%%%%%%%%%%%%%%%%%%%%%%%%%%%%%%%%%%%%%%%%%%%%%%%%
\subsection{Greatest Common Divisor}
\label{subsec:gcd-reference}
%%%%%%%%%%%%%%%%%%%%%%%%%%%%%%%%%%%%%%%%%%%%%%%%%%%%%%%%%%%%

The greatest common divisor is written

\[
\boxed{
\gcd(a,b).
}
\]

For example,

\[
\boxed{
\gcd(12,18)=6.
}
\]

%%%%%%%%%%%%%%%%%%%%%%%%%%%%%%%%%%%%%%%%%%%%%%%%%%%%%%%%%%%%
\subsection{Geometry Symbols}
\label{subsec:geometry-symbols-reference}
%%%%%%%%%%%%%%%%%%%%%%%%%%%%%%%%%%%%%%%%%%%%%%%%%%%%%%%%%%%%

\begin{center}
\begin{tabular}{c|l}
\textbf{Symbol} & \textbf{Meaning} \\ \hline
\( \angle ABC \) & angle \(ABC\) \\
\( \triangle ABC \) & triangle \(ABC\) \\
\( \perp \) & perpendicular \\
\( \parallel \) & parallel \\
\( \cong \) & congruent \\
\( \sim \) & similar, depending on context \\
\( \circ \) & degrees when used as a superscript
\end{tabular}
\end{center}

%%%%%%%%%%%%%%%%%%%%%%%%%%%%%%%%%%%%%%%%%%%%%%%%%%%%%%%%%%%%
\subsection{Trigonometric Symbols}
\label{subsec:trigonometric-symbols-reference}
%%%%%%%%%%%%%%%%%%%%%%%%%%%%%%%%%%%%%%%%%%%%%%%%%%%%%%%%%%%%

Common trigonometric functions are

\[
\boxed{
\sin x,
\quad
\cos x,
\quad
\tan x,
}
\]

\[
\boxed{
\csc x,
\quad
\sec x,
\quad
\cot x.
}
\]

Inverse trigonometric functions are commonly written

\[
\boxed{
\arcsin x,
\quad
\arccos x,
\quad
\arctan x.
}
\]

%%%%%%%%%%%%%%%%%%%%%%%%%%%%%%%%%%%%%%%%%%%%%%%%%%%%%%%%%%%%
\subsection{Topology Symbols}
\label{subsec:topology-symbols-reference}
%%%%%%%%%%%%%%%%%%%%%%%%%%%%%%%%%%%%%%%%%%%%%%%%%%%%%%%%%%%%

Common topology notation includes

\[
\boxed{
\operatorname{int}(A)
}
\]

for interior,

\[
\boxed{
\overline{A}
}
\]

for closure,

\[
\boxed{
\partial A
}
\]

for boundary, and

\[
\boxed{
A'
}
\]

for the derived set in some texts.

%%%%%%%%%%%%%%%%%%%%%%%%%%%%%%%%%%%%%%%%%%%%%%%%%%%%%%%%%%%%
\subsection{Neighborhoods}
\label{subsec:neighborhood-symbols-reference}
%%%%%%%%%%%%%%%%%%%%%%%%%%%%%%%%%%%%%%%%%%%%%%%%%%%%%%%%%%%%

An open ball of radius \(r\) centered at \(x\) is

\[
\boxed{
B_r(x)
=
\{
y:d(x,y)<r
\}.
}
\]

A closed ball is

\[
\boxed{
\overline{B}_r(x)
=
\{
y:d(x,y)\leq r
\}.
}
\]

%%%%%%%%%%%%%%%%%%%%%%%%%%%%%%%%%%%%%%%%%%%%%%%%%%%%%%%%%%%%
\subsection{Abstract Algebra Symbols}
\label{subsec:abstract-algebra-symbols-reference}
%%%%%%%%%%%%%%%%%%%%%%%%%%%%%%%%%%%%%%%%%%%%%%%%%%%%%%%%%%%%

Common algebraic structures are written:

\[
\boxed{
(G,\cdot)
}
\]

for a group,

\[
\boxed{
(R,+,\cdot)
}
\]

for a ring,

and

\[
\boxed{
(F,+,\cdot)
}
\]

for a field.

Substructures may be written

\[
\boxed{
H\leq G
}
\]

for subgroup and

\[
\boxed{
N\trianglelefteq G
}
\]

for normal subgroup.

%%%%%%%%%%%%%%%%%%%%%%%%%%%%%%%%%%%%%%%%%%%%%%%%%%%%%%%%%%%%
\subsection{Generated Subgroups}
\label{subsec:generated-subgroup-reference}
%%%%%%%%%%%%%%%%%%%%%%%%%%%%%%%%%%%%%%%%%%%%%%%%%%%%%%%%%%%%

The subgroup generated by \(g\) is

\[
\boxed{
\langle g\rangle.
}
\]

More generally,

\[
\boxed{
\langle S\rangle
}
\]

denotes the subgroup generated by a set \(S\).

%%%%%%%%%%%%%%%%%%%%%%%%%%%%%%%%%%%%%%%%%%%%%%%%%%%%%%%%%%%%
\subsection{Quotient Structures}
\label{subsec:quotient-symbols-reference}
%%%%%%%%%%%%%%%%%%%%%%%%%%%%%%%%%%%%%%%%%%%%%%%%%%%%%%%%%%%%

If \(N\trianglelefteq G\), then

\[
\boxed{
G/N
}
\]

denotes the quotient group.

For a ring \(R\) and ideal \(I\),

\[
\boxed{
R/I
}
\]

denotes the quotient ring.

%%%%%%%%%%%%%%%%%%%%%%%%%%%%%%%%%%%%%%%%%%%%%%%%%%%%%%%%%%%%
\subsection{Isomorphism Symbols}
\label{subsec:isomorphism-symbols-reference}
%%%%%%%%%%%%%%%%%%%%%%%%%%%%%%%%%%%%%%%%%%%%%%%%%%%%%%%%%%%%

The notation

\[
\boxed{
A\cong B
}
\]

often means \(A\) and \(B\) are isomorphic.

An isomorphism may be indicated by

\[
\boxed{
A\xrightarrow{\sim}B.
}
\]

%%%%%%%%%%%%%%%%%%%%%%%%%%%%%%%%%%%%%%%%%%%%%%%%%%%%%%%%%%%%
\subsection{Direct Sum}
\label{subsec:direct-sum-reference}
%%%%%%%%%%%%%%%%%%%%%%%%%%%%%%%%%%%%%%%%%%%%%%%%%%%%%%%%%%%%

The direct sum symbol is

\[
\boxed{
\oplus.
}
\]

For example,

\[
\boxed{
V
=
V_1\oplus V_2.
}
\]

It appears in linear algebra, group theory, representation theory, and
homological algebra.

%%%%%%%%%%%%%%%%%%%%%%%%%%%%%%%%%%%%%%%%%%%%%%%%%%%%%%%%%%%%
\subsection{Tensor Product}
\label{subsec:tensor-product-reference}
%%%%%%%%%%%%%%%%%%%%%%%%%%%%%%%%%%%%%%%%%%%%%%%%%%%%%%%%%%%%

The tensor product is written

\[
\boxed{
V\otimes W.
}
\]

The symbol

\[
\otimes
\]

also appears in tensor algebra, quantum mechanics, and multilinear algebra.

%%%%%%%%%%%%%%%%%%%%%%%%%%%%%%%%%%%%%%%%%%%%%%%%%%%%%%%%%%%%
\subsection{Graph Theory Symbols}
\label{subsec:graph-symbols-reference}
%%%%%%%%%%%%%%%%%%%%%%%%%%%%%%%%%%%%%%%%%%%%%%%%%%%%%%%%%%%%

A graph is

\[
\boxed{
G=(V,E).
}
\]

Common notation includes

\[
\boxed{
|V|
}
\]

for number of vertices,

\[
\boxed{
|E|
}
\]

for number of edges,

\[
\boxed{
\deg(v)
}
\]

for degree of vertex \(v\), and

\[
\boxed{
N(v)
}
\]

for its neighborhood.

%%%%%%%%%%%%%%%%%%%%%%%%%%%%%%%%%%%%%%%%%%%%%%%%%%%%%%%%%%%%
\subsection{Adjacency and Laplacian Matrices}
\label{subsec:graph-matrix-symbols-reference}
%%%%%%%%%%%%%%%%%%%%%%%%%%%%%%%%%%%%%%%%%%%%%%%%%%%%%%%%%%%%

The adjacency matrix is commonly

\[
\boxed{
A.
}
\]

The degree matrix is

\[
\boxed{
D.
}
\]

The graph Laplacian is

\[
\boxed{
L=D-A.
}
\]

%%%%%%%%%%%%%%%%%%%%%%%%%%%%%%%%%%%%%%%%%%%%%%%%%%%%%%%%%%%%
\subsection{Arrows}
\label{subsec:mathematical-arrows-reference}
%%%%%%%%%%%%%%%%%%%%%%%%%%%%%%%%%%%%%%%%%%%%%%%%%%%%%%%%%%%%

\begin{center}
\begin{tabular}{c|l}
\textbf{Symbol} & \textbf{Common Meaning} \\ \hline
\( \to \) & approaches, maps to, or tends to \\
\( \mapsto \) & maps an element to its image \\
\( \Rightarrow \) & implies \\
\( \Leftarrow \) & implied by \\
\( \Leftrightarrow \) & equivalent to \\
\( \longrightarrow \) & mapping or limiting process \\
\( \hookrightarrow \) & injection or embedding \\
\( \twoheadrightarrow \) & surjection \\
\( \xrightarrow{\sim} \) & isomorphism
\end{tabular}
\end{center}

%%%%%%%%%%%%%%%%%%%%%%%%%%%%%%%%%%%%%%%%%%%%%%%%%%%%%%%%%%%%
\subsection{Convergence Symbols}
\label{subsec:convergence-symbols-reference}
%%%%%%%%%%%%%%%%%%%%%%%%%%%%%%%%%%%%%%%%%%%%%%%%%%%%%%%%%%%%

Ordinary convergence is often written

\[
\boxed{
x_n\to x.
}
\]

Convergence in probability may be written

\[
\boxed{
X_n
\xrightarrow{P}
X.
}
\]

Convergence in distribution may be written

\[
\boxed{
X_n
\xrightarrow{d}
X.
}
\]

Almost-sure convergence may be written

\[
\boxed{
X_n
\xrightarrow{\mathrm{a.s.}}
X.
}
\]

%%%%%%%%%%%%%%%%%%%%%%%%%%%%%%%%%%%%%%%%%%%%%%%%%%%%%%%%%%%%
\subsection{Supremum and Infimum}
\label{subsec:supremum-infimum-reference}
%%%%%%%%%%%%%%%%%%%%%%%%%%%%%%%%%%%%%%%%%%%%%%%%%%%%%%%%%%%%

The least upper bound is

\[
\boxed{
\sup A.
}
\]

The greatest lower bound is

\[
\boxed{
\inf A.
}
\]

For a function,

\[
\boxed{
\sup_{x\in X}f(x)
}
\]

and

\[
\boxed{
\inf_{x\in X}f(x)
}
\]

need not be attained.

%%%%%%%%%%%%%%%%%%%%%%%%%%%%%%%%%%%%%%%%%%%%%%%%%%%%%%%%%%%%
\subsection{Maximum and Minimum}
\label{subsec:max-min-reference}
%%%%%%%%%%%%%%%%%%%%%%%%%%%%%%%%%%%%%%%%%%%%%%%%%%%%%%%%%%%%

A maximum is written

\[
\boxed{
\max_{x\in X}f(x),
}
\]

while a minimum is

\[
\boxed{
\min_{x\in X}f(x).
}
\]

Unlike supremum and infimum, maximum and minimum imply attainment.

%%%%%%%%%%%%%%%%%%%%%%%%%%%%%%%%%%%%%%%%%%%%%%%%%%%%%%%%%%%%
\subsection{Floor and Ceiling}
\label{subsec:floor-ceiling-reference}
%%%%%%%%%%%%%%%%%%%%%%%%%%%%%%%%%%%%%%%%%%%%%%%%%%%%%%%%%%%%

The floor is

\[
\boxed{
\lfloor x\rfloor,
}
\]

the greatest integer less than or equal to \(x\).

The ceiling is

\[
\boxed{
\lceil x\rceil,
}
\]

the least integer greater than or equal to \(x\).

%%%%%%%%%%%%%%%%%%%%%%%%%%%%%%%%%%%%%%%%%%%%%%%%%%%%%%%%%%%%
\subsection{Kronecker Delta}
\label{subsec:kronecker-delta-reference}
%%%%%%%%%%%%%%%%%%%%%%%%%%%%%%%%%%%%%%%%%%%%%%%%%%%%%%%%%%%%

The Kronecker delta is

\[
\boxed{
\delta_{ij}
=
\begin{cases}
1, & i=j,\\
0, & i\neq j.
\end{cases}
}
\]

It acts as a discrete identity tensor.

%%%%%%%%%%%%%%%%%%%%%%%%%%%%%%%%%%%%%%%%%%%%%%%%%%%%%%%%%%%%
\subsection{Dirac Delta}
\label{subsec:dirac-delta-reference}
%%%%%%%%%%%%%%%%%%%%%%%%%%%%%%%%%%%%%%%%%%%%%%%%%%%%%%%%%%%%

The Dirac delta is written

\[
\boxed{
\delta(x).
}
\]

It is a generalized function satisfying formally

\[
\boxed{
\int_{-\infty}^{\infty}
\delta(x-a)f(x)\,dx
=
f(a).
}
\]

It should not be confused with the Kronecker delta.

%%%%%%%%%%%%%%%%%%%%%%%%%%%%%%%%%%%%%%%%%%%%%%%%%%%%%%%%%%%%
\subsection{Inner Products}
\label{subsec:inner-product-reference}
%%%%%%%%%%%%%%%%%%%%%%%%%%%%%%%%%%%%%%%%%%%%%%%%%%%%%%%%%%%%

An abstract inner product is often written

\[
\boxed{
\langle x,y\rangle.
}
\]

In Euclidean space,

\[
\boxed{
\langle x,y\rangle
=
x^Ty.
}
\]

In function spaces,

\[
\boxed{
\langle f,g\rangle
=
\int
f(x)\overline{g(x)}\,dx
}
\]

under common conventions.

%%%%%%%%%%%%%%%%%%%%%%%%%%%%%%%%%%%%%%%%%%%%%%%%%%%%%%%%%%%%
\subsection{Indicator Functions}
\label{subsec:indicator-function-reference}
%%%%%%%%%%%%%%%%%%%%%%%%%%%%%%%%%%%%%%%%%%%%%%%%%%%%%%%%%%%%

An indicator function may be written

\[
\boxed{
\mathbf{1}_A(x)
}
\]

or

\[
\boxed{
\mathbb{I}_A(x).
}
\]

It is defined by

\[
\boxed{
\mathbf{1}_A(x)
=
\begin{cases}
1, & x\in A,\\
0, & x\notin A.
\end{cases}
}
\]

%%%%%%%%%%%%%%%%%%%%%%%%%%%%%%%%%%%%%%%%%%%%%%%%%%%%%%%%%%%%
\subsection{Characteristic Functions in Probability}
\label{subsec:probability-characteristic-function-reference}
%%%%%%%%%%%%%%%%%%%%%%%%%%%%%%%%%%%%%%%%%%%%%%%%%%%%%%%%%%%%

The characteristic function of random variable \(X\) is

\[
\boxed{
\varphi_X(t)
=
\mathbb{E}[e^{itX}].
}
\]

This use of ``characteristic function'' differs from the older use of the
term for an indicator function.

%%%%%%%%%%%%%%%%%%%%%%%%%%%%%%%%%%%%%%%%%%%%%%%%%%%%%%%%%%%%
\subsection{Common Greek Letters in Mathematics}
\label{subsec:greek-symbol-preview}
%%%%%%%%%%%%%%%%%%%%%%%%%%%%%%%%%%%%%%%%%%%%%%%%%%%%%%%%%%%%

Greek letters are widely used for parameters, constants, angles, and
variables.

Examples include:

\[
\boxed{
\alpha,\beta,\gamma,\delta,\epsilon,\theta,\lambda,\mu,\sigma,\phi,\omega.
}
\]

Capital forms include

\[
\boxed{
\Gamma,\Delta,\Theta,\Lambda,\Sigma,\Phi,\Omega.
}
\]

A complete Greek alphabet reference appears in the next section.

%%%%%%%%%%%%%%%%%%%%%%%%%%%%%%%%%%%%%%%%%%%%%%%%%%%%%%%%%%%%
\subsection{Commonly Confused Symbols}
\label{subsec:commonly-confused-symbols}
%%%%%%%%%%%%%%%%%%%%%%%%%%%%%%%%%%%%%%%%%%%%%%%%%%%%%%%%%%%%

Several symbols look similar but represent different objects.

For example:

\[
\boxed{
0
\neq
O
}
\]

where \(0\) is zero and \(O\) is the letter \(O\).

Similarly,

\[
\boxed{
1
\neq
l
\neq
I
}
\]

where the symbols may represent the number one, lowercase \(l\), or capital
\(I\).

%%%%%%%%%%%%%%%%%%%%%%%%%%%%%%%%%%%%%%%%%%%%%%%%%%%%%%%%%%%%
\subsection{Epsilon and Membership}
\label{subsec:epsilon-membership-confusion}
%%%%%%%%%%%%%%%%%%%%%%%%%%%%%%%%%%%%%%%%%%%%%%%%%%%%%%%%%%%%

The symbol

\[
\boxed{
\epsilon
}
\]

or

\[
\boxed{
\varepsilon
}
\]

is the Greek letter epsilon.

The symbol

\[
\boxed{
\in
}
\]

means ``is an element of.''

They are different symbols.

%%%%%%%%%%%%%%%%%%%%%%%%%%%%%%%%%%%%%%%%%%%%%%%%%%%%%%%%%%%%
\subsection{Subset and Less Than}
\label{subsec:subset-less-than-confusion}
%%%%%%%%%%%%%%%%%%%%%%%%%%%%%%%%%%%%%%%%%%%%%%%%%%%%%%%%%%%%

Do not confuse

\[
\boxed{
<
}
\]

with

\[
\boxed{
\subset.
}
\]

The first compares ordered quantities.

The second compares sets.

%%%%%%%%%%%%%%%%%%%%%%%%%%%%%%%%%%%%%%%%%%%%%%%%%%%%%%%%%%%%
\subsection{Element and Subset}
\label{subsec:element-subset-confusion}
%%%%%%%%%%%%%%%%%%%%%%%%%%%%%%%%%%%%%%%%%%%%%%%%%%%%%%%%%%%%

The notation

\[
\boxed{
x\in A
}
\]

compares an element with a set.

The notation

\[
\boxed{
A\subseteq B
}
\]

compares two sets.

For example,

\[
\boxed{
2\in\{1,2,3\},
}
\]

but

\[
\boxed{
\{2\}\subseteq\{1,2,3\}.
}
\]

%%%%%%%%%%%%%%%%%%%%%%%%%%%%%%%%%%%%%%%%%%%%%%%%%%%%%%%%%%%%
\subsection{Minus and Set Difference}
\label{subsec:minus-set-difference}
%%%%%%%%%%%%%%%%%%%%%%%%%%%%%%%%%%%%%%%%%%%%%%%%%%%%%%%%%%%%

For numbers,

\[
a-b
\]

usually means subtraction.

For sets,

\[
A\setminus B
\]

means set difference.

Using the dedicated symbol \(\setminus\) helps avoid ambiguity.

%%%%%%%%%%%%%%%%%%%%%%%%%%%%%%%%%%%%%%%%%%%%%%%%%%%%%%%%%%%%
\subsection{Vertical Bar Meanings}
\label{subsec:vertical-bar-meanings}
%%%%%%%%%%%%%%%%%%%%%%%%%%%%%%%%%%%%%%%%%%%%%%%%%%%%%%%%%%%%

The symbol

\[
|
\]

has several meanings.

Examples include:

\[
\boxed{
|x|
}
\]

for absolute value,

\[
\boxed{
|A|
}
\]

for cardinality,

\[
\boxed{
a\mid b
}
\]

for divisibility,

and

\[
\boxed{
\Pr(A\mid B)
}
\]

for conditional probability.

Context determines the meaning.

%%%%%%%%%%%%%%%%%%%%%%%%%%%%%%%%%%%%%%%%%%%%%%%%%%%%%%%%%%%%
\subsection{Tilde Meanings}
\label{subsec:tilde-meanings}
%%%%%%%%%%%%%%%%%%%%%%%%%%%%%%%%%%%%%%%%%%%%%%%%%%%%%%%%%%%%

The symbol

\[
\sim
\]

has many uses.

Examples include:

\[
\boxed{
X\sim N(0,1)
}
\]

for distribution,

\[
\boxed{
f(x)\sim g(x)
}
\]

for asymptotic equivalence,

and sometimes

\[
\boxed{
A\sim B
}
\]

for an equivalence relation or similarity.

The meaning must be inferred from context.

%%%%%%%%%%%%%%%%%%%%%%%%%%%%%%%%%%%%%%%%%%%%%%%%%%%%%%%%%%%%
\subsection{Star and Asterisk Notation}
\label{subsec:star-notation-reference}
%%%%%%%%%%%%%%%%%%%%%%%%%%%%%%%%%%%%%%%%%%%%%%%%%%%%%%%%%%%%

A superscript star may denote:

\[
\boxed{
x^*
}
\]

an optimal value,

\[
\boxed{
V^*
}
\]

a dual space,

or other context-dependent objects.

An asterisk has no single universal mathematical meaning.

%%%%%%%%%%%%%%%%%%%%%%%%%%%%%%%%%%%%%%%%%%%%%%%%%%%%%%%%%%%%
\subsection{Prime Notation}
\label{subsec:prime-notation-reference}
%%%%%%%%%%%%%%%%%%%%%%%%%%%%%%%%%%%%%%%%%%%%%%%%%%%%%%%%%%%%

A prime mark can denote:

\[
\boxed{
f'
}
\]

a derivative,

\[
\boxed{
x'
}
\]

a modified variable,

or

\[
\boxed{
A'
}
\]

a derived or transformed object.

Context is essential.

%%%%%%%%%%%%%%%%%%%%%%%%%%%%%%%%%%%%%%%%%%%%%%%%%%%%%%%%%%%%
\subsection{Superscript Minus One}
\label{subsec:superscript-minus-one-reference}
%%%%%%%%%%%%%%%%%%%%%%%%%%%%%%%%%%%%%%%%%%%%%%%%%%%%%%%%%%%%

The notation

\[
\boxed{
A^{-1}
}
\]

can mean matrix inverse,

while

\[
\boxed{
f^{-1}
}
\]

can mean inverse function.

Neither automatically means reciprocal.

For scalar \(a\neq0\),

\[
\boxed{
a^{-1}
=
\frac1a.
}
\]

%%%%%%%%%%%%%%%%%%%%%%%%%%%%%%%%%%%%%%%%%%%%%%%%%%%%%%%%%%%%
\subsection{Indices}
\label{subsec:index-notation-reference}
%%%%%%%%%%%%%%%%%%%%%%%%%%%%%%%%%%%%%%%%%%%%%%%%%%%%%%%%%%%%

Subscripts often label entries:

\[
\boxed{
x_i.
}
\]

Multiple indices may describe matrices or tensors:

\[
\boxed{
A_{ij},
\qquad
T_{ijk}.
}
\]

Index ranges should be stated when they are not obvious.

%%%%%%%%%%%%%%%%%%%%%%%%%%%%%%%%%%%%%%%%%%%%%%%%%%%%%%%%%%%%
\subsection{Einstein Summation Convention}
\label{subsec:einstein-summation-reference}
%%%%%%%%%%%%%%%%%%%%%%%%%%%%%%%%%%%%%%%%%%%%%%%%%%%%%%%%%%%%

In tensor notation, repeated indices may imply summation.

For example,

\[
\boxed{
a_ib_i
}
\]

may mean

\[
\boxed{
\sum_i a_ib_i.
}
\]

This is called the Einstein summation convention.

It should be announced before use.

%%%%%%%%%%%%%%%%%%%%%%%%%%%%%%%%%%%%%%%%%%%%%%%%%%%%%%%%%%%%
\subsection{Colon Product}
\label{subsec:colon-product-reference}
%%%%%%%%%%%%%%%%%%%%%%%%%%%%%%%%%%%%%%%%%%%%%%%%%%%%%%%%%%%%

For second-order tensors,

\[
\boxed{
A:B
=
\sum_{i,j}
A_{ij}B_{ij}.
}
\]

This is called a double contraction.

It appears frequently in continuum mechanics.

%%%%%%%%%%%%%%%%%%%%%%%%%%%%%%%%%%%%%%%%%%%%%%%%%%%%%%%%%%%%
\subsection{Differential Operators}
\label{subsec:differential-operators-reference}
%%%%%%%%%%%%%%%%%%%%%%%%%%%%%%%%%%%%%%%%%%%%%%%%%%%%%%%%%%%%

Common operators include

\[
\boxed{
\nabla
}
\]

for gradient-related operations,

\[
\boxed{
\Delta
}
\]

for the Laplacian,

\[
\boxed{
\frac{d}{dx}
}
\]

for ordinary differentiation,

and

\[
\boxed{
\frac{\partial}{\partial x}
}
\]

for partial differentiation.

%%%%%%%%%%%%%%%%%%%%%%%%%%%%%%%%%%%%%%%%%%%%%%%%%%%%%%%%%%%%
\subsection{Operator Names}
\label{subsec:operator-names-reference}
%%%%%%%%%%%%%%%%%%%%%%%%%%%%%%%%%%%%%%%%%%%%%%%%%%%%%%%%%%%%

Named operators should usually be typeset upright:

\[
\boxed{
\sin,
\cos,
\ln,
\exp,
\det,
\operatorname{tr},
\ker,
\operatorname{rank},
\sup,
\inf.
}
\]

This visually distinguishes operators from ordinary variables.

%%%%%%%%%%%%%%%%%%%%%%%%%%%%%%%%%%%%%%%%%%%%%%%%%%%%%%%%%%%%
\subsection{Kernel and Image}
\label{subsec:kernel-image-reference}
%%%%%%%%%%%%%%%%%%%%%%%%%%%%%%%%%%%%%%%%%%%%%%%%%%%%%%%%%%%%

For linear map

\[
T:V\to W,
\]

the kernel is

\[
\boxed{
\ker T
=
\{
v\in V:T(v)=0
\}.
}
\]

The image is

\[
\boxed{
\operatorname{im}T
=
\{
T(v):v\in V
\}.
}
\]

%%%%%%%%%%%%%%%%%%%%%%%%%%%%%%%%%%%%%%%%%%%%%%%%%%%%%%%%%%%%
\subsection{Dimension}
\label{subsec:dimension-reference}
%%%%%%%%%%%%%%%%%%%%%%%%%%%%%%%%%%%%%%%%%%%%%%%%%%%%%%%%%%%%

For vector space \(V\),

\[
\boxed{
\dim V
}
\]

denotes its dimension.

For example,

\[
\boxed{
\dim\R^n=n.
}
\]

%%%%%%%%%%%%%%%%%%%%%%%%%%%%%%%%%%%%%%%%%%%%%%%%%%%%%%%%%%%%
\subsection{Closure Under Context}
\label{subsec:notation-context-principle}
%%%%%%%%%%%%%%%%%%%%%%%%%%%%%%%%%%%%%%%%%%%%%%%%%%%%%%%%%%%%

No mathematical symbol has meaning completely independent of context.

For example,

\[
\boxed{
D
}
\]

may represent:

\[
\begin{itemize}
    \item a domain;
    \item a derivative operator;
    \item a degree matrix;
    \item a diffusion coefficient;
    \item a random variable;
    \item a determinant-related object.
\end{itemize}

Therefore symbols should be explicitly defined whenever ambiguity is
possible.

%%%%%%%%%%%%%%%%%%%%%%%%%%%%%%%%%%%%%%%%%%%%%%%%%%%%%%%%%%%%
\subsection{Recommended Notation Practices}
\label{subsec:recommended-notation-practices}
%%%%%%%%%%%%%%%%%%%%%%%%%%%%%%%%%%%%%%%%%%%%%%%%%%%%%%%%%%%%

Good mathematical writing should:

\begin{enumerate}
    \item define notation before using it;
    \item avoid changing meanings without warning;
    \item avoid using visually similar symbols for different objects;
    \item use boldface consistently for vectors when that convention is
          chosen;
    \item use uppercase and lowercase symbols systematically;
    \item distinguish scalars, vectors, matrices, sets, and operators;
    \item specify index ranges;
    \item specify domains and codomains of functions;
    \item state whether inequalities are strict or non-strict;
    \item state conventions that vary between texts;
    \item use notation to clarify mathematics rather than merely shorten it.
\end{enumerate}

%%%%%%%%%%%%%%%%%%%%%%%%%%%%%%%%%%%%%%%%%%%%%%%%%%%%%%%%%%%%
\subsection{Quick Symbol Reference}
\label{subsec:quick-symbol-reference}
%%%%%%%%%%%%%%%%%%%%%%%%%%%%%%%%%%%%%%%%%%%%%%%%%%%%%%%%%%%%

\begin{center}
\begin{tabular}{c|l}
\textbf{Symbol} & \textbf{Common Meaning} \\ \hline
\( = \) & equal \\
\( \neq \) & not equal \\
\( <,> \) & less than, greater than \\
\( \leq,\geq \) & less/greater than or equal \\
\( \approx \) & approximately equal \\
\( \in \) & element of \\
\( \notin \) & not an element of \\
\( \subseteq \) & subset of \\
\( \cup \) & union \\
\( \cap \) & intersection \\
\( \varnothing \) & empty set \\
\( \forall \) & for all \\
\( \exists \) & there exists \\
\( \Rightarrow \) & implies \\
\( \Leftrightarrow \) & if and only if \\
\( \sum \) & summation \\
\( \prod \) & product \\
\( \int \) & integral \\
\( \partial \) & partial derivative or boundary \\
\( \nabla \) & gradient/del operator \\
\( \Delta \) & Laplacian or change, depending on context \\
\( \infty \) & infinity \\
\( \propto \) & proportional to \\
\( \perp \) & perpendicular \\
\( \parallel \) & parallel \\
\( \sim \) & context-dependent equivalence/distribution/asymptotics \\
\( \cong \) & congruent or isomorphic, depending on context \\
\( \oplus \) & direct sum \\
\( \otimes \) & tensor product \\
\( \det \) & determinant \\
\( \operatorname{tr} \) & trace \\
\( \ker \) & kernel \\
\( \dim \) & dimension
\end{tabular}
\end{center}

%%%%%%%%%%%%%%%%%%%%%%%%%%%%%%%%%%%%%%%%%%%%%%%%%%%%%%%%%%%%
\subsection{Conceptual Summary}
\label{subsec:mathematical-symbols-summary}
%%%%%%%%%%%%%%%%%%%%%%%%%%%%%%%%%%%%%%%%%%%%%%%%%%%%%%%%%%%%

Mathematical notation provides a compact language for expressing mathematical
structures and relationships.

The major families of notation include:

\[
\boxed{
\text{arithmetic notation},
}
\]

\[
\boxed{
\text{set notation},
}
\]

\[
\boxed{
\text{logical notation},
}
\]

\[
\boxed{
\text{function notation},
}
\]

\[
\boxed{
\text{calculus notation},
}
\]

\[
\boxed{
\text{linear-algebra notation},
}
\]

\[
\boxed{
\text{probability notation},
}
\]

\[
\boxed{
\text{optimization notation},
}
\]

and

\[
\boxed{
\text{structural notation from advanced mathematics}.
}
\]

The same symbol can occasionally have different meanings in different
subjects.

Therefore the most important notation principle is

\[
\boxed{
\text{define symbols clearly and use them consistently}.
}
\]

%%%%%%%%%%%%%%%%%%%%%%%%%%%%%%%%%%%%%%%%%%%%%%%%%%%%%%%%%%%%
\subsection{Section Connections}
\label{subsec:mathematical-symbols-connections}
%%%%%%%%%%%%%%%%%%%%%%%%%%%%%%%%%%%%%%%%%%%%%%%%%%%%%%%%%%%%

\chapterconnection{
Mathematical Symbols and Notation begins the Mathematical Reference chapter
by collecting the symbolic language used throughout the textbook. Many of
these symbols were introduced gradually in Foundations and then reused in
Algebra, Number Theory, Geometry, Topology, Analysis, Probability,
Statistics, Differential Equations, Optimization, Numerical Mathematics,
Logic, and Applied Mathematics. The remaining sections of this chapter
organize specific reference material into focused tables and formula
collections. The next section gives the complete Greek alphabet together
with common mathematical uses of Greek letters.
}

%%%%%%%%%%%%%%%%%%%%%%%%%%%%%%%%%%%%%%%%%%%%%%%%%%%%%%%%%%%%
% SECTION SOURCE TRACKING
%%%%%%%%%%%%%%%%%%%%%%%%%%%%%%%%%%%%%%%%%%%%%%%%%%%%%%%%%%%%

%------------------------------------------------------------
% 14.1 Mathematical Symbols and Notation
%------------------------------------------------------------
%
% Source basis:
%   - Standard mathematical notation and conventions.
%   - Consolidated notation used throughout this textbook.
%
% Used for:
%   - arithmetic symbols;
%   - equality and inequalities;
%   - sets and membership;
%   - intervals;
%   - logic and quantifiers;
%   - functions and mappings;
%   - sequences;
%   - sums and products;
%   - calculus;
%   - vectors and matrices;
%   - complex numbers;
%   - probability;
%   - statistics;
%   - optimization;
%   - differential equations;
%   - asymptotic notation;
%   - number theory;
%   - geometry;
%   - topology;
%   - abstract algebra;
%   - graph theory;
%   - arrows and convergence;
%   - indices and tensors;
%   - operator notation;
%   - common notation ambiguities;
%   - notation-writing conventions.
%
% Notes:
%   - Greek Alphabet is developed next in Section 14.2.
%   - Number Systems follows in Section 14.3.
%
%%%%%%%%%%%%%%%%%%%%%%%%%%%%%%%%%%%%%%%%%%%%%%%%%%%%%%%%%%%%